\scanpage{78}
% Archive pencil foliation: 77.
\unsure{\deleted{d}D|\deleted{D}d|D|d}'autre part, on a un torseur évident $P_{\unsure{o|0}}$ sur $\handwritten{\underline}{Q}_p$ de groupe $(\handwritten{\hat}{\handwritten{\underline}{Z}})_{\inserted{\underline{Q}_p}}$, correspondant à l'e\unsure{xt|\deleted{x}t|\deleted{t}x} ension non ramifiée ma\unsure{\deleted{i}x|\deleted{x}i}imale de $\handwritten{\underline}{Q}_p$ ; \deleted{d}on\deleted{t} prend l'image \inserted{$P$}\handwritten{\begin{tikzpicture}[baseline=0cm,x=1cm,y=1cm]\draw[thin] (0,-.24)..controls(.08,0)and(.23,.16)..(.29,.53);\end{tikzpicture}} de ce torseur: c'est lui.

8.3. Soit $\handwritten{\underline}{H}$ la gerbe noyau de (8.2.1), qui a donc comme lien le groupe diagonalisable $H$ de 8.1. D'après le sor\unsure{\deleted{t}i|\deleted{i}t|i|t}\unsure{\deleted{i}t|\deleted{t}i|t|i}e général (ref) la classe
\margin{\handwritten{$\underline{H}$ est aussi la \uline{gerbe des relèvements}, \uline{pour l'hom.}
\diagram{\begin{tikzcd}[ampersand replacement=\&,column sep=small,row sep=small]
j:\ \underline{G}(F_p) \arrow[r] \arrow[d,phantom,"\simeq"{rotate=90}] \& \underline{\Pi}_p(F_p) \arrow[d,phantom,"\simeq"{rotate=90}] \\
[G] \& [G']
\end{tikzcd}}
déduit de l'hom \begin{tikzcd}[cramped]\unsure{G|\tilde{G}}\arrow[r]&G'\end{tikzcd}, du torseur $P$ ci-dessus.}}
\[
\handwritten{\xi\in}H^{\handwritten{2}}(\handwritten{\underline}{Q}_p,\deleted{\underline{G}_m)}\ H)=\operatorname{Hom}(\handwritten{\underline}{Q},\operatorname{Br}(\handwritten{\underline}{Q}_p))=\operatorname{Hom}(\handwritten{\underline}{Q},\handwritten{\underline}{Q}/\handwritten{\underline}{Z})
\]
de cette gerbe est donnée par
\[
\handwritten{\xi=\unsure{\partial|\delta}(\unsure{\eta|\gamma})}
\]
où
\[
\handwritten{\unsure{\eta|\gamma}\in}H^{\handwritten{1}}(\handwritten{\underline}{Q}_p,G')
\]
est la classe du torseur $P$. Pour fai\unsure{\deleted{e}r|\deleted{r}e|r}e le calcul, notons qu'on a un \deleted{sous-groupe}\inserted{\unsure{\deleted{e}q|\deleted{q}e|\deleted{a}q|\deleted{q}a|q}\unsure{\deleted{d}u|\deleted{u}d|\deleted{r}u|\deleted{u}r|u}otient} $H'$ de $G_r$, correspondant au sous-groupe de $\deleted{\unsure{Q_p|\underline{Q}_p}}$ $\handwritten{\overline}{\handwritten{\underline}{Q}}_p^{\handwritten{*}}$ \deleted{formée} engendré par les raci\unsure{\deleted{n}e|\deleted{e}n|e}\unsure{\deleted{e}n|\deleted{n}e|n}s de l'unité(= éléments de torsion) et toutes les raci\unsure{\deleted{n}e|\deleted{e}n|e}ns \unsure{\deleted{P}d|\deleted{d}P|\deleted{D}d|\deleted{d}D|d}e $p$, $p^{1/n}$ (ce groupe\deleted{\unsure{e|s}} est bien invariant par $\operatorname{Gal}(\handwritten{\overline}{\handwritten{\underline}{Q}}_p/\handwritten{\underline}{Q}_p)$ ). On a alors un homomorphisme d'extensions
\margin{\handwritten{\uline{NB} Ce diagramme est induit par un diagramme de \uline{gerbes} (valable sur un corps $k$ parfait quelconque) avec $\underline{H}'=$ gerbe correspondant à la sous-catégorie de $\operatorname{F\,Ciso}(\mathbb{F}_p)$ \unsure{formée|formé|formés} des \unsure{objets|objts} qui deviennent isomorphes sur une extension \uline{finie} $k'$ de $k=\mathbb{F}_p$ à une somme d'\unsure{objets|objts} $E_{r,s}$ ; et $H$ s'identifie à la \uline{gerbe des relèvements} de $P_{\unsure{o|0}}$ pour l'hom. \begin{tikzcd}[cramped]j:\underline{H}'\arrow[r]&[\hat{\underline{Z}}]\end{tikzcd} associé à \begin{tikzcd}[cramped]H'\arrow[r]&(\hat{\underline{Z}})_{Q_p}\end{tikzcd} ...}}
\diagram{\begin{tikzcd}[ampersand replacement=\&,execute at end picture={\draw[gray,thin] ([xshift=4pt,yshift=0pt]\tikzcdmatrixname-2-4.center)--([xshift=4pt,yshift=11pt]\tikzcdmatrixname-2-4.center);\draw[gray,thin] ([xshift=2pt,yshift=4pt]\tikzcdmatrixname-2-4.center)--([xshift=6pt,yshift=4pt]\tikzcdmatrixname-2-4.center)..controls([xshift=8pt,yshift=4pt]\tikzcdmatrixname-2-4.center)and([xshift=9pt,yshift=3pt]\tikzcdmatrixname-2-4.center)..([xshift=9pt,yshift=2pt]\tikzcdmatrixname-2-4.center);\fill[gray] ([xshift=9pt,yshift=-1pt]\tikzcdmatrixname-2-4.center) circle(.4pt);}]
0 \handwritten{\arrow[r]} \& H \handwritten{\arrow[r]} \handwritten{\arrow[d]} \& G \handwritten{\arrow[r,no head]} \handwritten{\arrow[d]} \& G' \handwritten{\arrow[r]} \handwritten{\arrow[d]} \& 0 \\
0 \handwritten{\arrow[r]} \& H \handwritten{\arrow[r]} \& H' \handwritten{\arrow[r]} \& \handwritten{(\hat{\underline{Z}})_{\unsure{\mathbb{Q}|Q}_p}} \handwritten{\arrow[r]} \& 0\ ,
\end{tikzcd}}
qui montre qu'\unsure{\deleted{i}o|\deleted{o}i}\unsure{\deleted{l}n|\deleted{n}l|\deleted{d}n|\deleted{n}d|n} \deleted{suffit de calculer} a
\[
\handwritten{\xi=\unsure{\partial|\delta}(\unsure{\eta|\gamma}_{\unsure{o|0}})},
\]
où
\[
\handwritten{\rlap{\begin{tikzpicture}[baseline=0cm,x=1cm,y=1cm]\draw[gray,thin] (.02,.24)..controls(-.02,.37)and(.13,.49)..(.12,.37)..controls(.09,.24)and(.01,.25)..(.02,.24);\draw[gray,thin] (.12,.29)..controls(.17,.46)and(.27,.47)..(.24,.36)--(.15,.24);\end{tikzpicture}}\unsure{\eta|\gamma}_{\unsure{o|0}}\in}H^{\handwritten{1}}(\handwritten{\underline}{Q}_p,\handwritten{\hat}{\handwritten{\underline}{Z}})
\]
est l'im\unsure{\deleted{d}a|\deleted{a}d|a}ge de $\handwritten{\unsure{\eta|\gamma}}$ par \begin{tikzcd}[cramped]G'\arrow[r]&\handwritten{\hat}{\handwritten{\underline}{Z}}\end{tikzcd}. On constate que le composé de (8.2.2) avec l'homomorphisme précédent est l'identité de $\handwritten{\hat}{\handwritten{\underline}{Z}}$, donc $\handwritten{\rlap{\begin{tikzpicture}[baseline=0cm,x=1cm,y=1cm]\draw[gray,thin] (.02,-.15)..controls(.2,-.09)and(.14,-.32)..(.03,-.24)..controls(-.01,-.2)and(.01,-.17)..(.02,-.15);\end{tikzpicture}}\unsure{\eta|\gamma}_{\unsure{o|0}}}$ est la classe du torseur canonique $P_{\unsure{o|0}}$ de la fin de 8.2. A l'aide\deleted{d} de ceci, il faut vérifier que la classe $\handwritten{\xi}$ correspond (par l'isomorphisme
\[
\operatorname{Br}(\handwritten{\underline}{Q}_p)\handwritten{\simeq}\handwritten{\underline}{Q}_{\deleted{\handwritten{p}}}/\handwritten{\underline}{Z}_{\deleted{\handwritten{p}}}
\]
du corps de classe \unsure{\deleted{c}l|\deleted{l}c}ocal, à l'homomorphisme canonique \begin{tikzcd}[cramped]\handwritten{\mathbb{Q}}\arrow[r]&\handwritten{\mathbb{Q}/\mathbb{Z}}\end{tikzcd}

